% This document was generated by an LLM.

\documentclass{essay}

\title{L'H\^opital's Rule: Statement and Proof}
\author{}
\date{}

\begin{document}

\maketitle

\section{Prerequisite}

Rolle's theorem is taken as known and is not reproved here.

\begin{theorem}[Rolle]
  Let $h:[p,q]\to\mathbb{R}$ be continuous on $[p,q]$ and differentiable on $(p,q)$, with $h(p)=h(q)$. Then there exists $\xi\in(p,q)$ such that $h'(\xi)=0$.
\end{theorem}

\section{Cauchy's Mean Value Theorem}

\begin{lemma}[Cauchy]
  Let $f,g:[p,q]\to\mathbb{R}$ be continuous on $[p,q]$ and differentiable on $(p,q)$. Then there exists $\xi\in(p,q)$ such that
  \[
    [f(q)-f(p)]\,g'(\xi) \;=\; [g(q)-g(p)]\,f'(\xi).
  \]
  If in addition $g'(x)\neq 0$ for every $x\in(p,q)$, then $g(p)\neq g(q)$ and
  \[
    \frac{f(q)-f(p)}{g(q)-g(p)} \;=\; \frac{f'(\xi)}{g'(\xi)}.
  \]
\end{lemma}

\begin{proof}
  Define $h:[p,q]\to\mathbb{R}$ by
  \[
    h(x) = [f(q)-f(p)]\,g(x) - [g(q)-g(p)]\,f(x).
  \]
  As a linear combination of $f$ and $g$, $h$ is continuous on $[p,q]$ and differentiable on $(p,q)$. A direct computation gives
  \[
    h(p) = f(q)g(p) - g(q)f(p), \qquad h(q) = f(q)g(p) - g(q)f(p),
  \]
  so $h(p)=h(q)$. By Rolle's theorem there is $\xi\in(p,q)$ with $h'(\xi)=0$; since
  \[
    h'(x) = [f(q)-f(p)]\,g'(x) - [g(q)-g(p)]\,f'(x),
  \]
  this is exactly the claimed identity.

  Now suppose $g'(x)\neq 0$ throughout $(p,q)$. If $g(p)=g(q)$, Rolle's theorem applied to $g$ itself on $[p,q]$ would produce a point of $(p,q)$ at which $g'$ vanishes, contradicting the hypothesis. Hence $g(p)\neq g(q)$, and since also $g'(\xi)\neq 0$ (else the identity above would force $g(p)=g(q)$ as well, again a contradiction), dividing by $g'(\xi)[g(q)-g(p)]$ yields the second formula.
\end{proof}

\section{L'H\^opital's Rule}

\begin{theorem}[L'H\^opital]
  Let $-\infty \le c < d \le +\infty$, and let $f,g:(c,d)\to\mathbb{R}$ be differentiable, with $g'(x)\neq 0$ for every $x\in(c,d)$. Suppose
  \[
    \lim_{x\to c^+} \frac{f'(x)}{g'(x)} = A, \qquad A\in\mathbb{R}\cup\{-\infty,+\infty\},
  \]
  (when $c=-\infty$, ``$x\to c^+$'' is read simply as $x\to-\infty$), and suppose in addition that either
  \begin{enumerate}
    \item[(a)] $\displaystyle \lim_{x\to c^+} f(x) = \lim_{x\to c^+} g(x) = 0$, \quad or
    \item[(b)] $\displaystyle \lim_{x\to c^+} g(x) = +\infty$ \; (or $-\infty$).
  \end{enumerate}
  Then
  \[
    \lim_{x\to c^+} \frac{f(x)}{g(x)} = A.
  \]
\end{theorem}

\begin{proof}
  \textbf{Step 0: $g$ is injective on $(c,d)$.} If $g(x)=g(y)$ for some $x\neq y$ in $(c,d)$, Rolle's theorem applied to $g$ on the interval between $x$ and $y$ would produce a zero of $g'$ in $(c,d)$, contradicting the hypothesis $g'\neq 0$. So $g$ takes each value at most once on $(c,d)$; being also continuous, it is strictly monotonic there (a standard consequence of the intermediate value theorem).

  \medskip
  \textbf{Step 1: the case $A\in\mathbb{R}$.} Fix $\varepsilon>0$ and choose real numbers $q,r$ with
  \[
    A-\varepsilon < q < A < r < A+\varepsilon.
  \]
  Since $f'/g'\to A$ as $x\to c^+$, there is a point $b\in(c,d)$ such that
  \[
    q < \frac{f'(t)}{g'(t)} < r \qquad \text{for all } t\in(c,b). \tag{1}
  \]

  For any $x,y$ with $c<x<y<b$, Step 0 gives $g(x)\neq g(y)$, so Cauchy's theorem (Lemma 1) applies on $[x,y]$: there exists $\xi=\xi(x,y)\in(x,y)\subset(c,b)$ with
  \[
    \frac{f(x)-f(y)}{g(x)-g(y)} = \frac{f'(\xi)}{g'(\xi)}.
  \]
  Combined with (1), this gives
  \[
    q < \frac{f(x)-f(y)}{g(x)-g(y)} < r \qquad \text{for all } c<x<y<b. \tag{2}
  \]

  \emph{Case (a).} Fix $y\in(c,b)$. By Step 0, $g$ is strictly monotonic on $(c,d)$ with $\lim_{x\to c^+}g(x)=0$, so $0$ is not attained on $(c,d)$: hence $g(y)\neq 0$. Letting $x\to c^+$ in (2), and using $f(x)\to 0$, $g(x)\to 0$, the middle term tends to $f(y)/g(y)$, and weak inequalities survive the limit:
  \[
    q \;\le\; \frac{f(y)}{g(y)} \;\le\; r \qquad \text{for every } y\in(c,b). \tag{3}
  \]

  \emph{Case (b).} Fix $y\in(c,b)$. Because $g(x)\to+\infty$ as $x\to c^+$ (the case $g\to-\infty$ is identical after replacing $g$ by $-g$), we may choose $x$ close enough to $c$ that $g(x)>0$ and $g(x)>g(y)$. Multiplying (2) by the positive quantity $[g(x)-g(y)]/g(x)$ gives
  \[
    q\Bigl(1-\frac{g(y)}{g(x)}\Bigr) \;<\; \frac{f(x)}{g(x)} - \frac{f(y)}{g(x)} \;<\; r\Bigl(1-\frac{g(y)}{g(x)}\Bigr). \tag{4}
  \]
  Let $x\to c^+$ with $y$ fixed. Since $g(x)\to+\infty$, both $g(y)/g(x)\to 0$ and $f(y)/g(x)\to 0$, so (4) yields
  \[
    q \;\le\; \liminf_{x\to c^+}\frac{f(x)}{g(x)} \;\le\; \limsup_{x\to c^+}\frac{f(x)}{g(x)} \;\le\; r. \tag{5}
  \]

  In either case, (3) or (5) gives
  \[
    q \;\le\; \liminf_{x\to c^+}\frac{f(x)}{g(x)} \;\le\; \limsup_{x\to c^+}\frac{f(x)}{g(x)} \;\le\; r.
  \]
  This holds for \emph{every} admissible pair $q,r$ with $A-\varepsilon<q<A<r<A+\varepsilon$. Letting $\varepsilon\to 0^+$ forces $q\to A$, $r\to A$, hence
  \[
    \liminf_{x\to c^+}\frac{f(x)}{g(x)} = \limsup_{x\to c^+}\frac{f(x)}{g(x)} = A,
  \]
  i.e., $\displaystyle\lim_{x\to c^+} f(x)/g(x) = A$.

  \medskip
  \textbf{Step 2: the case $A=+\infty$ (the case $A=-\infty$ is symmetric).} Fix any $r\in\mathbb{R}$. Since $f'/g'\to+\infty$, there is $b\in(c,d)$ with $f'(t)/g'(t) > r$ for all $t\in(c,b)$. Repeating the argument of Step 1 using only this one-sided bound (there is no companion $q$ to track) produces, in either case (a) or (b),
  \[
    \liminf_{x\to c^+}\frac{f(x)}{g(x)} \;\ge\; r.
  \]
  Since $r\in\mathbb{R}$ was arbitrary, $\liminf_{x\to c^+} f(x)/g(x) = +\infty$, i.e., $\lim_{x\to c^+} f(x)/g(x) = +\infty$.
\end{proof}

\section{Extensions}

\begin{corollary}[Left-hand and two-sided limits]
  The theorem holds equally for $x\to c^-$: if $f,g$ are differentiable on $(e,c)$ for some $e<c$, with $g'\neq 0$ there, and $\lim_{x\to c^-} f'(x)/g'(x) = A$ together with the analogous hypothesis (a) or (b), then $\lim_{x\to c^-} f(x)/g(x) = A$. Indeed, set $\tilde f(x) = f(2c-x)$ and $\tilde g(x)=g(2c-x)$ on $(c,\,2c-e)$; as $x\to c^+$, $2c-x\to c^-$, and $\tilde f'(x)/\tilde g'(x) = f'(2c-x)/g'(2c-x) \to A$, so Theorem 2 applies to $\tilde f,\tilde g$ and yields the claim after undoing the substitution. Consequently, if $f,g$ are differentiable on a full punctured neighborhood $(e,c)\cup(c,d)$ of a finite point $c$ and both one-sided derivative-ratio limits equal $A$, then $\lim_{x\to c} f(x)/g(x) = A$.
\end{corollary}

\begin{corollary}[Limits at infinity]
  The case $x\to+\infty$ reduces to Theorem 2 via $t=1/x$: if $F(t)=f(1/t)$, $G(t)=g(1/t)$ for $t$ near $0^+$, the chain rule gives
  \[
    \frac{F'(t)}{G'(t)} = \frac{f'(1/t)}{g'(1/t)},
  \]
  so $\lim_{x\to+\infty} f'(x)/g'(x) = A$ implies $\lim_{t\to 0^+} F'(t)/G'(t) = A$; hypotheses (a)/(b) transfer identically, so Theorem 2 gives $\lim_{t\to 0^+} F(t)/G(t) = A$, which is $\lim_{x\to+\infty} f(x)/g(x) = A$. The case $x\to-\infty$ is analogous via $t=-1/x$.
\end{corollary}

\begin{remark}
  The hypothesis that $\lim_{x\to c^+} f'(x)/g'(x)$ \emph{exists} cannot be dropped, and the theorem is not reversible. Take $f(x) = x+\sin x$, $g(x)=x$ on $(0,\infty)$; both diverge to $+\infty$ as $x\to\infty$, and
  \[
    \frac{f(x)}{g(x)} = 1+\frac{\sin x}{x} \;\longrightarrow\; 1,
  \]
  so the quotient converges. Yet
  \[
    \frac{f'(x)}{g'(x)} = 1+\cos x
  \]
  oscillates between $0$ and $2$ and has no limit as $x\to\infty$. This is not a contradiction: the theorem states a \emph{sufficient}, not necessary, condition for evaluating $\lim f/g$.
\end{remark}

\end{document}
